\Needspace{9\baselineskip}\begin{center}\textit{Lösung der unbestimmten Gleichung \(mxx+nyy=A\) durch Ausschlüsse.}\end{center}


\Needspace{6\baselineskip}
\Article{323}
In Abschnitt V haben wir gezeigt, wie man sämtliche Darstellungen einer gegebenen Zahl \(A\) durch die binäre Form \(mxx+nyy\),
\SourceLeaf{397}
also die Lösungen der unbestimmten Gleichung \(mxx+nyy=A\), nach einem allgemeinen Verfahren findet, dessen Kürze ebenfalls nichts zu wünschen übrig zu lassen scheint, wenn sämtliche Werte des Ausdrucks \(\surd{-mn}\) nach dem Modul \(A\) selbst und nach der durch ihre quadratischen Faktoren geteilten Zahl bereits vorliegen. Hier wollen wir jedoch für den Fall, dass \(mn\) positiv ist, eine Lösung erläutern, die viel bequemer als die direkte ist, wenn man für diese jene Werte erst berechnen muss. Wir setzen dabei voraus, dass die Zahlen \(m\), \(n\) und \(A\) positiv und paarweise teilerfremd sind, da sich die übrigen Fälle leicht auf diesen zurückführen lassen. Offenbar genügt es auch, positive Werte von \(x,\allowbreak y\) zu ermitteln, weil die übrigen daraus allein durch Änderung der Vorzeichen folgen.

Offenbar muss \(x\) so beschaffen sein, dass \(\frac{A-mxx}{n}\), wofür wir \(V\) schreiben wollen, positiv, ganzzahlig und ein Quadrat wird. Die erste Bedingung verlangt, dass \(x\) nicht größer als \(\surd{\frac Am}\) ist; die zweite ist bei \(n=1\) von selbst erfüllt und verlangt andernfalls, dass der Wert des Ausdrucks \(\frac Am\; (\mathrm{mod}.\ n)\) ein quadratischer Rest von \(n\) ist. Bezeichnen wir sämtliche verschiedenen Werte des Ausdrucks \(\surd{\frac Am}\; (\mathrm{mod}.\ n)\) mit \(\pm r,\allowbreak \pm r'\) usw., so muss \(x\) unter einer der Formen \(nt+r,\allowbreak nt-r,\allowbreak nt+r'\) usw. enthalten sein. Am einfachsten wäre es daher, sämtliche Zahlen dieser Formen unterhalb der Grenze \(\surd{\frac Am}\) (die Gesamtheit dieser Zahlen wollen wir mit \(\Omega\) bezeichnen) für \(x\) einzusetzen und nur diejenigen beizubehalten, für die \(V\) ein Quadrat wird. Wie man diesen Versuch beliebig verkürzen kann, werden wir im nächsten Artikel zeigen.

\Needspace{6\baselineskip}
\Article{324}
Das Ausschlussverfahren, mit dem wir dies erreichen wollen, besteht wie in der vorigen Untersuchung darin, mehrere Zahlen nach Belieben zu wählen, die auch hier \textit{Ausschlusszahlen} heißen sollen, zu untersuchen, für welche Werte von \(x\) der Wert von \(V\) ein quadratischer Nichtrest dieser Ausschlusszahlen wird, und solche Werte von \(x\) aus \(\Omega\) zu entfernen. Durch Überlegungen, die den in Artikel \(321\) dargelegten völlig entsprechen, ergibt sich, dass man nur Primzahlen oder Primzahlpotenzen als Ausschlusszahlen verwenden muss. Bei einer Ausschlusszahl der letzteren Art sind unter den Werten von \(V\) nur diejenigen ihrer Nichtreste auszuschließen, die Reste sämtlicher niedrigerer Potenzen derselben Primzahl sind, sofern die Ausschlüsse mit diesen bereits vorgenommen worden sind.

Die Ausschlusszahl sei also \(E=p^\mu\), einschließlich des Falles \(\mu=1\), wobei \(p\) eine Primzahl ist, die \(m\) nicht teilt, und wir nehmen an\HistoricalNote{*)}{Der Kürze wegen behandeln wir die beiden Fälle, in denen \(n\) durch \(p\) teilbar beziehungsweise nicht teilbar ist, gemeinsam; im letzteren muss man \(\nu=0\) setzen.}, dass \(p^\nu\) die höchste Potenz
\SourceLeaf{398}
derselben Primzahl ist, durch die \(n\) teilbar ist. Es seien \(a,\allowbreak b,\allowbreak c\) usw. quadratische Nichtreste von \(E\): sämtliche, wenn \(\mu=1\), und die notwendigen, also diejenigen, die Reste der niedrigeren Potenzen sind, wenn \(\mu>1\). Man berechne die Wurzeln der Kongruenzen \(mz\equiv A-na,\allowbreak mz\equiv A-nb,\allowbreak mz\equiv A-nc\) usw. \((\mathrm{mod.}\ Ep^\nu=p^{\mu+\nu})\), die \(\alpha,\allowbreak \beta,\allowbreak \gamma\) usw. seien. Dann ist leicht einzusehen, dass, wenn für einen Wert von \(x\) die Kongruenz \(xx\equiv\alpha\) \((\mathrm{mod.}\ Ep^\nu)\) gilt, der entsprechende Wert von \(V\) \(\equiv a\; (\mathrm{mod}.\ E)\), also ein Nichtrest von \(E\) wird; ebenso verhält es sich mit den übrigen Zahlen \(\beta,\allowbreak \gamma\) usw. Umgekehrt ist ebenso leicht einzusehen, dass, wenn ein Wert von \(x\) die Kongruenz \(V\equiv a\; (\mathrm{mod}.\ E)\) liefert, für denselben Wert \(xx\equiv\alpha\; (\mathrm{mod}.\ Ep^\nu)\) gilt. Somit liefern alle Werte von \(x\), für die \(xx\) keiner der Zahlen \(\alpha,\allowbreak \beta,\allowbreak \gamma\) usw. nach dem Modul \(Ep^\nu\) kongruent ist, solche Werte von \(V\), die keiner der Zahlen \(a,\allowbreak b,\allowbreak c\) usw. nach dem Modul \(E\) kongruent sind. Man wähle nun unter den Zahlen \(\alpha,\allowbreak \beta,\allowbreak \gamma\) usw. sämtliche quadratischen Reste von \(Ep^\nu\) aus, die \(g,\allowbreak g',\allowbreak g''\) usw. seien, berechne die Werte der Ausdrücke \(\surd g,\allowbreak \surd{g'}.\surd{g''}\) usw. \((\mathrm{mod.}\ Ep^\nu)\) und nehme an, dass sich daraus \(\pm h,\allowbreak \pm h',\allowbreak \pm h''\) usw. ergeben. Dann ist klar, dass man sämtliche Zahlen der Formen \(Ep^\nu t\pm h,\allowbreak Ep^\nu t\pm h',\allowbreak Ep^\nu t\pm h''\) usw. sicher aus \(\Omega\) entfernen kann und dass keinem Wert von \(x\), der nach diesem Ausschluss in \(\Omega\) verbleibt, ein Wert von \(V\) der Formen \(Eu+a,\allowbreak Eu+b,\allowbreak Eu+c\) usw. entsprechen kann. Ferner ist offenbar, dass solche Werte von \(V\) schon von selbst aus keinem Wert von \(x\) entstehen können, wenn sich unter den Zahlen \(\alpha,\allowbreak \beta,\allowbreak \gamma\) usw. keine quadratischen Reste von \(Ep^\nu\) finden; in diesem Fall kann man die Zahl \(E\) daher nicht als Ausschlusszahl verwenden. — Man kann beliebig viele solche Ausschlusszahlen anwenden und so die Menge der Zahlen in \(\Omega\) beliebig verkleinern.

\enlargethispage{\baselineskip}
Sehen wir nun zu, ob man auch Primzahlen, die \(m\) teilen, oder Potenzen solcher Primzahlen als Ausschlusszahlen verwenden darf. Es sei \(B\) der Wert des Ausdrucks \(\frac An\; (\mathrm{mod}.\ m)\). Offenbar wird \(V\) stets \(B\) nach dem Modul \(m\) kongruent, welcher Wert auch für \(x\) genommen wird; zur Lösbarkeit der gegebenen Gleichung ist daher notwendig, dass \(B\) ein quadratischer Rest von \(m\) ist. Bezeichnet nun \(p\) einen beliebigen ungeraden Primteiler von \(m\), der nach Voraussetzung weder \(n\) noch \(A\) teilt und somit auch \(B\) nicht teilt, so ist für jeden Wert von \(x\) \(V\) ein Rest von \(p\) und damit auch jeder Potenz von \(p\); daher können \(p\) und seine Potenzen nicht als Ausschlusszahlen dienen. Auf genau dieselbe Weise ist, wenn \(m\) durch \(8\) teilbar ist, zur Lösbarkeit der gegebenen Gleichung notwendig, dass \(B\equiv1\; (\mathrm{mod}.\ 8)\) gilt. Dann wird auch \(V\) für jeden Wert von \(x\) \(\equiv1\) \((\mathrm{mod.}\ 8)\), und daher eignen sich die Zweierpotenzen nicht zum Ausschluss. — Wenn aber
\SourceLeaf{399}
\(m\) durch \(4\), jedoch nicht durch \(8\) teilbar ist, so muss aus demselben Grund \(B\equiv1\) \((\mathrm{mod.}\ 4)\) gelten; der Wert des Ausdrucks \(\frac An\; (\mathrm{mod}.\ 8)\) ist daher entweder \(1\) oder \(5\) und werde mit \(C\) bezeichnet. Man sieht leicht, dass für einen geraden Wert von \(x\) hier \(V\equiv C\), für einen ungeraden dagegen \(V\equiv C+4\; (\mathrm{mod}.\ 8)\) gilt. Daher sind die geraden Werte zu verwerfen, wenn \(C=5\), und die ungeraden, wenn \(C=1\). — Ist schließlich \(m\) durch \(2\), jedoch nicht durch \(4\) teilbar, so sei wie zuvor \(C\) der Wert des Ausdrucks \(\frac An\; (\mathrm{mod}.\ 8)\), der \(1,\allowbreak 3,\allowbreak 5\) oder \(7\) ist; und \(D\) sei der Wert von \(\frac{\frac12m}{n}\; (\mathrm{mod}.\ 4)\), der \(1\) oder \(3\) ist. Da nun der Wert von \(V\) offenbar stets \(\equiv C-2Dxx\; (\mathrm{mod}.\ 8)\) wird, also für gerades \(x\) \(\equiv C\) und für ungerades \(\equiv C-2D\), folgt daraus leicht, dass alle ungeraden Werte von \(x\) zu verwerfen sind, wenn \(C=1\); alle geraden, wenn \(C=3\) und \(D=1\) oder \(C=7\) und \(D=3\); und dass sämtliche übrigen Werte \(V\equiv1\; (\mathrm{mod}.\ 8)\), also einen Rest jeder Zweierpotenz liefern. In den anderen Fällen dagegen, nämlich wenn \(C=5\) oder \(C=3\) und \(D=3\) oder \(C=7\) und \(D=1\), wird \(V\equiv3,\allowbreak 5\) oder \(7\; (\mathrm{mod}.\ 8)\), ob man \(x\) gerade oder ungerade nimmt. Daraus geht hervor, dass die gegebene Gleichung in diesen Fällen überhaupt keine Lösung besitzt.

Da wir im Übrigen auf genau dieselbe Weise, wie wir hier den Wert von \(x\) durch Ausschlüsse zu finden gezeigt haben, mit den nötigen Änderungen auch den Wert von \(y\) ermitteln können, lässt sich das Ausschlussverfahren stets auf zwei Arten zur Lösung des gegebenen Problems anwenden, sofern nicht \(m=n=1\), wobei beide zusammenfallen. Meist ist diejenige vorzuziehen, für die \(\Omega\) weniger Glieder enthält, was man leicht im Voraus abschätzen kann. — Schließlich brauchen wir kaum zu bemerken, dass, wenn nach einigen Ausschlüssen \textit{sämtliche} Zahlen aus \(\Omega\) verschwunden sind, dies als sicheres Zeichen der Unlösbarkeit der gegebenen Gleichung anzusehen ist.

\Needspace{6\baselineskip}
\Article{325}
\textit{Beispiel.} Es sei die Gleichung \(3xx+455yy=10857362\) gegeben, die wir auf zwei Arten lösen wollen: \textit{erstens}, indem wir die Werte von \(x\), dann die Werte von \(y\) ermitteln. Die Grenze der ersteren ist in diesem Fall \(\surd{3619120\frac23}\), die zwischen \(1902\) und \(1903\) liegt; der Wert des Ausdrucks \(\frac A3\; (\mathrm{mod}.\ 455)\) ist \(354\), und die Werte des Ausdrucks \(\surd{354}\; (\mathrm{mod}.\ 455)\) sind \(\pm82,\allowbreak \pm152\), \(\pm173,\allowbreak \pm212\). Daher besteht \(\Omega\) aus den folgenden \(33\) Zahlen: \(82,\allowbreak 152,\allowbreak 173,\allowbreak 212,\allowbreak \) \(243,\allowbreak 282,\allowbreak 303,\allowbreak 373,\allowbreak 537,\allowbreak 607,\allowbreak 628,\allowbreak 667,\allowbreak 698,\allowbreak 737,\allowbreak 758,\allowbreak 828,\allowbreak 992,\allowbreak 1062,\allowbreak 1083,\allowbreak \) \(1122,\allowbreak 1153,\allowbreak 1192,\allowbreak 1213,\allowbreak 1283,\allowbreak 1447,\allowbreak 1517,\allowbreak 1538,\allowbreak 1577,\allowbreak 1608,\allowbreak 1647,\allowbreak 1668,\allowbreak 1738,\allowbreak \) \(1902\). Die Zahl \(3\) lässt sich in diesem Fall nicht zum Ausschluss verwenden, weil sie \(m\) teilt. Für die Ausschlusszahl \(4\) haben wir \(a=2,\allowbreak b=3\), also \(\alpha=0,\allowbreak \beta=3\);
\SourceLeaf{400}
\(g=0\), und die Werte des Ausdrucks \(\surd{g}\ (\mathrm{mod.}\ 4)\) sind \(0\) und \(2\). Somit sind alle Zahlen der Formen \(4t\) und \(4t+2\), \textit{d. h.} alle geraden Zahlen, aus \(\Omega\) zu entfernen; die sechzehn übrigen seien mit \(\Omega'\) bezeichnet. Für \(E=5\), wobei diese Ausschlusszahl auch \(n\) teilt, haben die Kongruenzen \(mz\equiv A-2n\) und \(mz\equiv A-3n\ (\mathrm{mod.}\ 25)\) die Wurzeln \(9\) und \(24\), die beide Reste von \(25\) sind; die Werte der Ausdrücke \(\surd9\) und \(\surd{24}\ (\mathrm{mod.}\ 25)\) sind \(\pm3,\allowbreak \ \pm7\). Entfernt man aus \(\Omega'\) sämtliche Zahlen der Formen \(25t\pm3,\allowbreak \ 25t\pm7\), so bleiben diese zehn \((\Omega'')\): \(173,\allowbreak 373,\allowbreak 537,\allowbreak 667,\allowbreak 737,\allowbreak 1083,\allowbreak 1213,\allowbreak 1283,\allowbreak 1517,\allowbreak 1577\). Für \(E=7\) haben die Kongruenzen \(mz\equiv A-3n,\allowbreak \ mz\equiv A-5n,\allowbreak \ mz\equiv A-6n\) \((\mathrm{mod.}\ 49)\) die Wurzeln \(32,\allowbreak 39,\allowbreak 18\), die sämtlich Reste von \(49\) sind, und die Werte der Ausdrücke \(\surd{32},\allowbreak \surd{39},\allowbreak \surd{18}\ (\mathrm{mod.}\ 49)\) sind \(\pm9,\allowbreak \pm23,\allowbreak \pm19\). Entfernt man aus \(\Omega''\) die Zahlen der Formen \(49t\pm9,\allowbreak \ 49t\pm19,\allowbreak \ 49t\pm23\), so bleiben diese fünf \((\Omega''')\): \(537,\allowbreak 737\), \(1083,\allowbreak 1213,\allowbreak 1517\). Für \(E=8\) haben wir \(a=5\), also \(\alpha=5\), einen Nichtrest von \(8\); daher lässt sich die Ausschlusszahl \(8\) nicht verwenden. Die Zahl \(9\) ist aus demselben Grund zu überspringen wie \(3\). Für \(E=11\) sind die Zahlen \(a,\allowbreak b\) usw. gleich \(2,\allowbreak 6\), \(7,\allowbreak 8,\allowbreak 10\); \(\nu=0\); somit \(\alpha,\allowbreak \beta\) usw. \(=8,\allowbreak 10,\allowbreak 5,\allowbreak 0,\allowbreak 1\), von denen drei Reste von \(11\) sind, nämlich \(0,\allowbreak 1,\allowbreak 5\). Daraus folgt, dass aus \(\Omega'''\) die Zahlen der Formen \(11t,\allowbreak \ 11t\pm1,\allowbreak \ 11t\pm4\) zu verwerfen sind, worauf \(537,\allowbreak 1083,\allowbreak 1213\) übrig bleiben. Bei der Probe ergeben sich für \(V\) der Reihe nach die Werte \(21961,\allowbreak 16129,\allowbreak 14161\), von denen allein der zweite und dritte Quadrate sind. Die gegebene Gleichung besitzt daher zwei Lösungen mit positiven Werten von \(x,\allowbreak y\), nämlich \(x=1083,\allowbreak y=127\) und \(x=1213,\allowbreak y=119\).

\begin{samepage}
\textit{Zweitens.} Will man die andere Unbekannte derselben Gleichung durch Ausschlüsse ermitteln, so bringe man diese in die Form \(455xx+3yy=10857362\), indem man \(x\) mit \(y\) vertauscht, um sämtliche Bezeichnungen der Artikel \(323,\allowbreak 324\) beibehalten zu können. Die Grenze der Werte von \(x\) liegt hier zwischen \(154\) und \(155\); der Wert des Ausdrucks \(\frac Am\ (\mathrm{mod.}\ n)\) ist \(1\), und die Werte von \(\surd1\ (\mathrm{mod.}\ 3)\) sind \(+1\) und \(-1\). Daher enthält \(\Omega\) sämtliche Zahlen der Formen \(3t+1\) und \(3t-1\), \textit{d. h.} alle durch \(3\) nicht teilbaren Zahlen bis einschließlich \(154\), deren Anzahl \(103\) ist. Wendet man die oben gegebenen Regeln an, so findet man, dass für die Ausschlusszahlen \(3;\allowbreak \ 4;\allowbreak \ 9;\allowbreak \ 11;\allowbreak \) \(17;\allowbreak \ 19;\allowbreak \ 23\) die Zahlen der Formen \(9t\pm4;\allowbreak \ 4t,\allowbreak \ 4t\pm2\), also alle geraden Zahlen; \(27t\pm1,\allowbreak \ 27t\pm10;\allowbreak \ 11t,\allowbreak \ 11t\pm1,\allowbreak \ 11t\pm3;\allowbreak \ 17t\pm3,\allowbreak \ 17t\pm4,\allowbreak \ 17t\pm5\), \(17t\pm7;\allowbreak \ 19t\pm2,\allowbreak \ 19t\pm3,\allowbreak \ 19t\pm8,\allowbreak \ 19t\pm9;\allowbreak \ 23t,\allowbreak \ 23t\pm1,\allowbreak \ 23t\pm5,\allowbreak \ 23t\pm7\), \(23t\pm9,\allowbreak \ 23t\pm10\) zu verwerfen sind. Nach deren Entfernung bleiben \(119,\allowbreak 127\) übrig, die allein \(V\) einen quadratischen Wert geben und dieselben Lösungen liefern, zu denen wir oben gelangt sind.
\par\end{samepage}
\SourceLeaf{401}
